\documentclass[]{article}

% Pacotes básicos
\usepackage{fontspec}
\usepackage{amsmath,amssymb}
\usepackage{unicode-math}

% Configuração de fontes
\setmainfont{Liberation Serif}
\setmathfont{Liberation Serif}

% Cores
\usepackage[dvipsnames,svgnames,table]{xcolor}
\definecolor{fondpaille}{cmyk}{0, 0.0118, 0.1176, 0}
\pagecolor{fondpaille}

% Línguas
\usepackage{polyglossia}
\setdefaultlanguage{portuguese}
\setotherlanguages{english}

% Layout - AUMENTEI a margem superior para o cabeçalho
\usepackage[a4paper,left=2.5cm,right=2.5cm,top=3.5cm,bottom=2.5cm]{geometry}
\linespread{1.3}

% Gráficos
\usepackage{graphicx}
\usepackage{tikz}
\usetikzlibrary{calc, positioning, arrows}

% Tabelas
\usepackage{booktabs}

% Configuração de cabeçalhos e rodapés
\usepackage{fancyhdr}
\pagestyle{fancy}
\fancyhf{} % Limpa cabeçalho e rodapé
\renewcommand{\headrulewidth}{0pt} % Remove linha do cabeçalho
\renewcommand{\footrulewidth}{0pt} % Remove linha do rodapé

% Remove cabeçalho da primeira página de cada capítulo
\fancypagestyle{plain}{
  \fancyhf{}
  \renewcommand{\headrulewidth}{0pt}
}

% Citações
\usepackage{csquotes}
\usepackage[
  backend=biber,
  style=numeric,
  language=auto
]{biblatex}

% ---------------------------------------------------------
% Topo e rodapé de páginas (logos e texto) - CORRIGIDO
% ---------------------------------------------------------
\usepackage{eso-pic}

\AddToShipoutPictureBG{%
  \ifnum\value{page}>2
    \AtPageUpperLeft{%
      \begin{tikzpicture}[remember picture, overlay]
        % Ajustei as posições para ficarem mais para cima
        \draw[line width=0.5pt, color=black] (2.5cm, -2.5cm) -- (\paperwidth - 2.5cm, -2.5cm);
        \draw[line width=0.5pt, color=black] (2.5cm, -\paperheight + 2.5cm) -- (\paperwidth - 2.5cm, -\paperheight + 2.5cm);

        % Conteúdo página ímpar - posições ajustadas
        \ifodd\value{page}
          \node [anchor=west, text depth=0pt, font=\tiny] at (2.5cm, -2.2cm) {\href{mailto:13132@aluno.aebenfica.pt}{Grupo Disciplinar 510}};
          % Logo menor e mais para cima
          \node [anchor=north east] at (\paperwidth-2.5cm, -1.8cm) {\includegraphics[width=0.04\paperwidth]{logos/aeB05.pdf}};
          \node [anchor=west, text height=0pt, font=\tiny] at (2.5cm, -\paperheight + 2.2cm) {\href{mailto:15803@aluno.aebenfica.pt}{Adrian Dias}};
          \node [anchor=west, font=\tiny] at (2.5cm, -\paperheight + 2.0cm) {\jobname.tex \smash{\rule[-0.5ex]{0.6pt}{1ex}}};
        % Conteúdo página par - posições ajustadas
        \else
          \node [anchor=west, text depth=0pt] at (2.5cm, -2.2cm) {\includegraphics[width=0.04\paperwidth]{logos/aeB05.pdf}};
          \node [anchor=east, text depth=0pt, font=\tiny] at (\paperwidth - 2.5cm, -2.2cm) {\href{mailto:13137@aluno.aebenfica.pt}{Grupo Disciplinar 510}};
          \node [anchor=east, text height=0pt, font=\tiny] at (\paperwidth - 2.5cm, -\paperheight + 2.2cm) {\href{mailto:13137@aluno.aebenfica.pt}{Adrian Dias}};
          \node [anchor=east, font=\tiny] at (\paperwidth - 2.5cm, -\paperheight + 2.0cm) {\jobname.tex \smash{\rule[-0.5ex]{0.6pt}{1ex}}};
        \fi
      \end{tikzpicture}
    }
  \fi
}

% hyperref (último)
\usepackage[colorlinks=true, linkcolor=black, urlcolor=blue]{hyperref}
